\documentclass[11pt,a4paper]{article}

\usepackage[T1]{fontenc}
\usepackage[utf8]{inputenc}
\usepackage[margin=2.5cm]{geometry}
\usepackage{booktabs}
\usepackage{hyperref}
\usepackage{microtype}
\usepackage{xcolor}

\hypersetup{hidelinks}
\newcommand{\code}[1]{\texttt{#1}}

\title{Issue Report Classification\\\large NLBSE'24 Project Report}
\author{\textcolor{red}{Add student names and matriculation numbers}}
\date{\today}

\begin{document}
\maketitle

\begin{abstract}
This project implements a modular Python application for classifying GitHub
issue reports as bugs, features, or questions. It provides interchangeable
in-memory and file-backed repositories, a leakage-aware preprocessing and
cross-validation pipeline, four classical TF--IDF classifiers, and a shared
evaluation layer compatible with the supplied Sentence Transformer baselines.
The deep-learning track adds a feed-forward network, a convolutional text
classifier, and fine-tuned DistilBERT. All model selection is performed on the
official training data before evaluation on the official test split.
\end{abstract}

\section{Problem and dataset}
The NLBSE'24 dataset contains 3,000 labelled issues from React, TensorFlow,
VS Code, Bitcoin, and OpenCV, spanning March 2016 to September 2023. Each
repository contributes 100 examples of each class to both the training and
test split. Although the course description calls the task multi-label, the
released benchmark excludes issues with multiple labels; the implemented task
is therefore three-class, single-label classification.

The competition trains one classifier per repository. Its headline score is
the arithmetic mean of the five per-repository weighted F1 scores. Macro F1,
accuracy, per-class scores, and pooled metrics are retained as secondary
diagnostics.

\section{Application architecture}
The domain entity is \code{IssueReport}. Business logic depends on the
\code{IssueRepository} abstract interface rather than a storage technology.
\code{InMemoryIssueRepository} stores objects in a Python list, while
\code{FileIssueRepository} persists the same fields in CSV or JSON. Both
implementations expose the same loading, saving, filtering, transformation,
and tabular-conversion operations. Equivalence tests run the same workflow
through both implementations and verify identical results and lossless file
round trips.

The \code{IssueDataService} coordinates loading, validation, preprocessing,
auditing, and fold creation. This separation keeps persistence code out of the
classification and evaluation modules and satisfies the requirement that the
persistence layer can be changed with minimal application changes.

\section{Data preparation}
Validation checks the schema, allowed labels and repositories, missing text,
cell balance, and split size. Duplicate detection hashes Unicode-normalized,
case-folded title/body text. Stratified group folds keep every normalized
duplicate group on one side of a fold, preventing vocabulary and label leakage.

Four cleaning levels are available: raw whitespace normalization,
software-aware conservative replacement, light Markdown/code cleanup, and
aggressive full normalization. A training-only ablation also varies title
weight, word truncation, and character n-grams. The best tested configuration
uses raw normalized text, a title weight of three, a 400-word limit, and both
word and character TF--IDF.

\begin{table}[ht]
\centering
\caption{Top preprocessing configurations on duplicate-safe training folds.}
\begin{tabular}{lrrlr}
\toprule
Cleaning & Title weight & Max words & Char TF-IDF & F1 \\
\midrule
raw & 3 & 400 & yes & 0.7543 \\
raw & 1 & 400 & yes & 0.7500 \\
raw & 3 & unlimited & yes & 0.7483 \\
raw & 3 & 400 & no & 0.7471 \\
raw & 1 & 400 & no & 0.7471 \\
raw & 3 & 200 & no & 0.7458 \\
raw & 1 & unlimited & yes & 0.7450 \\
raw & 3 & 200 & yes & 0.7445 \\
\bottomrule
\end{tabular}

\end{table}

\section{Classical models and selection}
The four candidates are Complement Naive Bayes, logistic regression, linear
SVM, and random forest. Each estimator is a complete scikit-learn pipeline, so
TF--IDF vocabulary and inverse-document frequencies are learned only from the
current training fold. Word features use unigrams and bigrams; character
features use boundary-aware 3--5 grams.

Model families are compared with five-fold stratified group cross-validation
over five random seeds. Selection uses the mean cross-repository weighted F1.
The first seed's complete out-of-fold predictions and all repeated-run
summaries are saved as versioned JSON for inspection.

\begin{table}[ht]
\centering
\caption{Repeated grouped cross-validation used for model selection.}
\begin{tabular}{lrr}
\toprule
Model & Mean weighted F1 & SD \\
\midrule
Logistic regression & 0.7506 & 0.0061 \\
Linear SVM & 0.7472 & 0.0057 \\
Random forest & 0.7149 & 0.0026 \\
Complement Naive Bayes & 0.6929 & 0.0074 \\
\bottomrule
\end{tabular}

\end{table}

\section{Deep-learning models}
Track C implements three models through the same evaluator. The feed-forward
network consumes word unigram/bigram TF--IDF and contains hidden layers of 256
and 64 ReLU units. The TextCNN learns 128-dimensional word embeddings and uses
parallel convolution filters of widths 2--5 followed by global max pooling.
Both use dropout, AdamW, gradient clipping, and validation-loss early stopping.

The transformer alternative is \emph{distilbert-base-uncased}, selected instead
of RoBERTa because the specification requires only one of the two and
DistilBERT has lower memory and runtime requirements on Colab. The complete
encoder and a new three-class head are fine-tuned with linear warmup.

Each model makes a duplicate-safe stratified validation split from its current
training partition. This happens before fitting the FFNN vectorizer or CNN
vocabulary, avoiding representation leakage. The official test set is never
used for early stopping. Per-epoch training loss, validation loss, validation
accuracy, weighted F1, best epoch, device, and parameter count are retained in
the result JSON.

\IfFileExists{tables/deep_learning.tex}{
\begin{table}[ht]
\centering
\caption{Completed Track C evaluations on the official holdout.}
\input{tables/deep_learning.tex}
\end{table}
}{
Final Track C scores are intentionally omitted until the corresponding
experiment artifacts are generated on the target CPU or GPU environment.
}

\section{Evaluation and robustness}
The selected classical family is retrained on the complete training subset of
each repository and evaluated once on its matching official test subset. The shared
evaluator records metrics, confusion matrices, timings, stable issue IDs, and
probabilities when the estimator exposes them. A project-stratified bootstrap
quantifies uncertainty. Error artifacts summarize confusion directions,
performance by text-length quintile, and high-confidence mistakes.

For a temporal robustness check, the combined dataset is deduplicated before
each repository/class cell is divided into older training issues and newer test
issues. A class- and repository-matched random split is evaluated as a control.
This is an additional robustness analysis, not a replacement for the official
competition split.

\begin{table}[ht]
\centering
\caption{Selected model on the official holdout split.}
\begin{tabular}{lrrr}
\toprule
Repository & Accuracy & Macro F1 & Weighted F1 \\
\midrule
bitcoin/bitcoin & 0.7100 & 0.7049 & 0.7049 \\
facebook/react & 0.8000 & 0.7984 & 0.7984 \\
microsoft/vscode & 0.7533 & 0.7532 & 0.7532 \\
opencv/opencv & 0.7500 & 0.7478 & 0.7478 \\
tensorflow/tensorflow & 0.7700 & 0.7699 & 0.7699 \\
Cross-repository mean & 0.7567 & 0.7548 & 0.7548 \\
\bottomrule
\end{tabular}

\end{table}

The classical winner outperforms the supplied fastText baseline but remains
below RoBERTa and the dedicated SetFit Sentence Transformer baseline. This is
consistent with pretrained representations transferring semantic information
that sparse features must learn from only 300 training issues per repository.

\begin{table}[ht]
\centering
\caption{Comparison with the supplied locally reproduced baselines.}
\begin{tabular}{lr}
\toprule
Model & Cross-repository weighted F1 \\
\midrule
SetFit (Sentence Transformer) & 0.8240 \\
RoBERTa & 0.7923 \\
TF-IDF + Logistic Regression & 0.7548 \\
fastText & 0.7184 \\
\bottomrule
\end{tabular}

\end{table}

\begin{table}[ht]
\centering
\caption{Bootstrap interval and approximate power diagnostic.}
\begin{tabular}{lr}
\toprule
Quantity & Value \\
\midrule
Point estimate & 0.7548 \\
95 percent bootstrap lower & 0.7329 \\
95 percent bootstrap upper & 0.7763 \\
Approximate detectable accuracy difference & 0.0300 \\
\bottomrule
\end{tabular}

\end{table}

\begin{table}[ht]
\centering
\caption{Matched random and chronological robustness evaluation.}
\begin{tabular}{lr}
\toprule
Protocol & Weighted F1 \\
\midrule
Matched random split & 0.7757 \\
Chronological split & 0.6733 \\
Chronological minus random & -0.1024 \\
\bottomrule
\end{tabular}

\end{table}

\section{Reproducibility}
The repository provides a development container, pinned dependency groups,
Jupyter notebooks, command-line experiment drivers, unit tests, and GitHub
Actions checks. The principal commands are \code{make data},
\code{make pipeline}, \code{make classical-ablation}, \code{make classical},
\code{make results}, and \code{make report}. Generated JSON retains model
settings and random seeds, allowing tables to be rebuilt without copying
numbers manually into this document. Track C is run with
\code{make deep-neural} and \code{make deep-transformer}; a GPU is strongly
recommended for the transformer.

\section{Limitations and conclusion}
The benchmark is small and balanced, while real issue streams are larger,
imbalanced, and change over time. Cross-validation folds are not independent,
and the minimum-detectable-difference calculation is an approximate paired
accuracy diagnostic rather than a formal confidence bound for weighted F1.
The temporal split also changes the training distribution and is intended as a
sensitivity analysis. Within these limits, the project supplies a complete,
modular, leakage-aware comparison of classical, neural, and pretrained
transformer models, with an auditable path from raw data to reported results.

\begin{thebibliography}{9}
\bibitem{nlbse24}
Kallis et al. \emph{The NLBSE'24 Tool Competition on Issue Report
Classification}. NLBSE, 2024.

\bibitem{setfit}
Tunstall et al. \emph{Efficient Few-Shot Learning Without Prompts}. 2022.

\bibitem{distilbert}
Sanh et al. \emph{DistilBERT, a distilled version of BERT: smaller, faster,
cheaper and lighter}. 2019.
\end{thebibliography}

\end{document}
